% AUTO-GENERATED by the update_record tool. DO NOT EDIT.
% verify_paper regenerates this file from record.json and the run outputs
% and fails on any difference, so manual edits always surface.
\noindent\textbf{Hypothesis 1.} SciGym 公開コードの Controller を論文条件（eval\_debug\_rounds=3、temperature=1.0）のまま用い、open-weight 3 モデル（DeepSeek-V4.1-Flash、Qwen3.8-27B、Nemotron-3-Ultra）を自前 serving して SciGym-small 137 件と SciGym-large 213 件を解かせたとき、反復上限を論文の 20 から 80 に増やすと平均 Simulation Trajectory Error（STE）は雑音幅を超えて（0.03 以上）改善する。~\cite[s1.p3, s1.p5, s1.p6]{duan-2025-measuring}, \cite[s3.p1]{scigymrikyuopenmodels-2026-scigym}
\begin{enumerate}
\item[\textbf{Claim 1}] DeepSeek-V4.1-Flash で SciGym-small 137 件を解かせたとき、反復上限 80 の平均 STE は反復上限 20（論文条件）より 0.03 以上低い。~\cite[s1.p3, s1.p4, s1.p5]{duan-2025-measuring}, \cite[s2.p1]{h4duan-2025-scigym}
  \emph{Rationale:} DeepSeek-V4.1-Flash の small split で、論文条件を超える予算が雑音幅（137 件平均の標準誤差約 0.02、差で約 0.03）を超えて効くことを示し、h1 の当該モデル・split の部分を担う。
  \emph{Criterion:} \texttt{\detokenize{deepseek-v4.1-flash-small-i80}}.\texttt{\detokenize{trajectory_smape}} $-$ \texttt{\detokenize{deepseek-v4.1-flash-small-i20}}.\texttt{\detokenize{trajectory_smape}} $\leq -0.03$.
  \emph{Prediction:} $[-0.15, -0.02]$ (先行する 100 反復の内部試行では kimi-k3 が 20$\to$100 反復で STE 0.225$\to$0.088（平均 67.8 反復を使用）、qwen3.6-35b は平均 20.7 反復で自発提出し 0.457$\to$0.446。予算を使い切るモデルほど改善が大きいと見て、モデルごとに幅を置く).
  \emph{Observed:} -0.0643907.
  \emph{Verdict:} supported.
\item[\textbf{Claim 2}] DeepSeek-V4.1-Flash で SciGym-large 213 件を解かせたとき、反復上限 80 の平均 STE は反復上限 20（論文条件）より 0.03 以上低い。~\cite[s1.p3, s1.p4, s1.p5]{duan-2025-measuring}, \cite[s2.p1]{h4duan-2025-scigym}
  \emph{Rationale:} DeepSeek-V4.1-Flash の large split で、論文条件を超える予算が雑音幅（137 件平均の標準誤差約 0.02、差で約 0.03）を超えて効くことを示し、h1 の当該モデル・split の部分を担う。
  \emph{Criterion:} \texttt{\detokenize{deepseek-v4.1-flash-large-i80}}.\texttt{\detokenize{trajectory_smape}} $-$ \texttt{\detokenize{deepseek-v4.1-flash-large-i20}}.\texttt{\detokenize{trajectory_smape}} $\leq -0.03$.
  \emph{Prediction:} $[-0.15, -0.02]$ (先行する 100 反復の内部試行では kimi-k3 が 20$\to$100 反復で STE 0.225$\to$0.088（平均 67.8 反復を使用）、qwen3.6-35b は平均 20.7 反復で自発提出し 0.457$\to$0.446。予算を使い切るモデルほど改善が大きいと見て、モデルごとに幅を置く).
  \emph{Observed:} -0.1495.
  \emph{Verdict:} supported.
\item[\textbf{Claim 3}] Qwen3.8-27B で SciGym-small 137 件を解かせたとき、反復上限 80 の平均 STE は反復上限 20（論文条件）より 0.03 以上低い。~\cite[s1.p3, s1.p4, s1.p5]{duan-2025-measuring}, \cite[s2.p1]{h4duan-2025-scigym}
  \emph{Rationale:} Qwen3.8-27B の small split で、論文条件を超える予算が雑音幅（137 件平均の標準誤差約 0.02、差で約 0.03）を超えて効くことを示し、h1 の当該モデル・split の部分を担う。
  \emph{Criterion:} \texttt{\detokenize{qwen3.8-27b-small-i80}}.\texttt{\detokenize{trajectory_smape}} $-$ \texttt{\detokenize{qwen3.8-27b-small-i20}}.\texttt{\detokenize{trajectory_smape}} $\leq -0.03$.
  \emph{Prediction:} $[-0.08, 0.02]$ (先行する 100 反復の内部試行では kimi-k3 が 20$\to$100 反復で STE 0.225$\to$0.088（平均 67.8 反復を使用）、qwen3.6-35b は平均 20.7 反復で自発提出し 0.457$\to$0.446。予算を使い切るモデルほど改善が大きいと見て、モデルごとに幅を置く).
  \emph{Observed:} -0.0973588.
  \emph{Verdict:} supported.
\item[\textbf{Claim 4}] Qwen3.8-27B で SciGym-large 213 件を解かせたとき、反復上限 80 の平均 STE は反復上限 20（論文条件）より 0.03 以上低い。~\cite[s1.p3, s1.p4, s1.p5]{duan-2025-measuring}, \cite[s2.p1]{h4duan-2025-scigym}
  \emph{Rationale:} Qwen3.8-27B の large split で、論文条件を超える予算が雑音幅（137 件平均の標準誤差約 0.02、差で約 0.03）を超えて効くことを示し、h1 の当該モデル・split の部分を担う。
  \emph{Criterion:} \texttt{\detokenize{qwen3.8-27b-large-i80}}.\texttt{\detokenize{trajectory_smape}} $-$ \texttt{\detokenize{qwen3.8-27b-large-i20}}.\texttt{\detokenize{trajectory_smape}} $\leq -0.03$.
  \emph{Prediction:} $[-0.08, 0.02]$ (先行する 100 反復の内部試行では kimi-k3 が 20$\to$100 反復で STE 0.225$\to$0.088（平均 67.8 反復を使用）、qwen3.6-35b は平均 20.7 反復で自発提出し 0.457$\to$0.446。予算を使い切るモデルほど改善が大きいと見て、モデルごとに幅を置く).
  \emph{Observed:} -0.130822.
  \emph{Verdict:} supported.
\item[\textbf{Claim 5}] Nemotron-3-Ultra で SciGym-small 137 件を解かせたとき、反復上限 80 の平均 STE は反復上限 20（論文条件）より 0.03 以上低い。~\cite[s1.p3, s1.p4, s1.p5]{duan-2025-measuring}, \cite[s2.p1]{h4duan-2025-scigym}
  \emph{Rationale:} Nemotron-3-Ultra の small split で、論文条件を超える予算が雑音幅（137 件平均の標準誤差約 0.02、差で約 0.03）を超えて効くことを示し、h1 の当該モデル・split の部分を担う。
  \emph{Criterion:} \texttt{\detokenize{nemotron-3-ultra-small-i80}}.\texttt{\detokenize{trajectory_smape}} $-$ \texttt{\detokenize{nemotron-3-ultra-small-i20}}.\texttt{\detokenize{trajectory_smape}} $\leq -0.03$.
  \emph{Prediction:} $[-0.12, -0.02]$ (先行する 100 反復の内部試行では kimi-k3 が 20$\to$100 反復で STE 0.225$\to$0.088（平均 67.8 反復を使用）、qwen3.6-35b は平均 20.7 反復で自発提出し 0.457$\to$0.446。予算を使い切るモデルほど改善が大きいと見て、モデルごとに幅を置く).
  \emph{Observed:} -0.0280823.
  \emph{Verdict:} refuted.
\item[\textbf{Claim 6}] Nemotron-3-Ultra で SciGym-large 213 件を解かせたとき、反復上限 80 の平均 STE は反復上限 20（論文条件）より 0.03 以上低い。~\cite[s1.p3, s1.p4, s1.p5]{duan-2025-measuring}, \cite[s2.p1]{h4duan-2025-scigym}
  \emph{Rationale:} Nemotron-3-Ultra の large split で、論文条件を超える予算が雑音幅（137 件平均の標準誤差約 0.02、差で約 0.03）を超えて効くことを示し、h1 の当該モデル・split の部分を担う。
  \emph{Criterion:} \texttt{\detokenize{nemotron-3-ultra-large-i80}}.\texttt{\detokenize{trajectory_smape}} $-$ \texttt{\detokenize{nemotron-3-ultra-large-i20}}.\texttt{\detokenize{trajectory_smape}} $\leq -0.03$.
  \emph{Prediction:} $[-0.12, -0.02]$ (先行する 100 反復の内部試行では kimi-k3 が 20$\to$100 反復で STE 0.225$\to$0.088（平均 67.8 反復を使用）、qwen3.6-35b は平均 20.7 反復で自発提出し 0.457$\to$0.446。予算を使い切るモデルほど改善が大きいと見て、モデルごとに幅を置く).
  \emph{Observed:} -0.0411857.
  \emph{Verdict:} supported.
\end{enumerate}
\noindent\emph{Assumed, so that the claims together imply Hypothesis 1:}
\begin{itemize}
\item temperature 1.0 での 1 回実行の平均 STE の実行間変動が判定に対して十分小さいこと（137 件平均の標準誤差約 0.02）（c1〜c12）
\item 自前 serving（vLLM / NIM）したモデルが、提供元の推奨設定で呼んだ場合と同じ能力で応答すること。思考はサーバー側の reasoning parser で本文から分離する（c1〜c12）
\item 文脈長の拡張（YaRN、NIM\_MAX\_MODEL\_LEN）が 262,144 以内の応答品質を変えないこと（c1〜c12）
\item aarch64 向けの libroadrunner 2.7.0 と antimony 2.14.0 が、論文環境と同じ軌道と評価値を与えること（c1〜c12）
\end{itemize}
\noindent\textbf{Hypothesis 2.} 同じ条件で反復上限を 40 から 80 に倍増しても、平均 STE の改善は 0.05 未満にとどまる（論文条件の 2 倍程度で性能は飽和する）。~\cite[s1.p3]{duan-2025-measuring}
\begin{enumerate}
\item[\textbf{Claim 7}] DeepSeek-V4.1-Flash で SciGym-small 137 件を解かせたとき、反復上限 40 から 80 への平均 STE の改善は 0.05 未満である。~\cite[s1.p3, s1.p4]{duan-2025-measuring}, \cite[s2.p1]{h4duan-2025-scigym}
  \emph{Rationale:} DeepSeek-V4.1-Flash の small split で、予算を 40 から 80 に倍増しても伸びが判定幅に収まる（飽和している）ことを示し、h2 の当該モデル・split の部分を担う。基準 run は c1 の反復上限 80。
  \emph{Criterion:} \texttt{\detokenize{deepseek-v4.1-flash-small-i40}}.\texttt{\detokenize{trajectory_smape}} $-$ \texttt{\detokenize{deepseek-v4.1-flash-small-i80}}.\texttt{\detokenize{trajectory_smape}} $\leq 0.05$.
  \emph{Prediction:} $[0, 0.08]$ (内部試行では kimi-k3 が 100 反復でも平均 67.8 反復で提出しており、40 を超える予算の使用は部分的。qwen 系は 20 前後で自発提出する。40$\to$80 の差は雑音幅（約 0.03）程度と見る).
  \emph{Observed:} 0.0390427.
  \emph{Verdict:} supported.
\item[\textbf{Claim 8}] DeepSeek-V4.1-Flash で SciGym-large 213 件を解かせたとき、反復上限 40 から 80 への平均 STE の改善は 0.05 未満である。~\cite[s1.p3, s1.p4]{duan-2025-measuring}, \cite[s2.p1]{h4duan-2025-scigym}
  \emph{Rationale:} DeepSeek-V4.1-Flash の large split で、予算を 40 から 80 に倍増しても伸びが判定幅に収まる（飽和している）ことを示し、h2 の当該モデル・split の部分を担う。基準 run は c2 の反復上限 80。
  \emph{Criterion:} \texttt{\detokenize{deepseek-v4.1-flash-large-i40}}.\texttt{\detokenize{trajectory_smape}} $-$ \texttt{\detokenize{deepseek-v4.1-flash-large-i80}}.\texttt{\detokenize{trajectory_smape}} $\leq 0.05$.
  \emph{Prediction:} $[0, 0.08]$ (内部試行では kimi-k3 が 100 反復でも平均 67.8 反復で提出しており、40 を超える予算の使用は部分的。qwen 系は 20 前後で自発提出する。40$\to$80 の差は雑音幅（約 0.03）程度と見る).
  \emph{Observed:} 0.0583697.
  \emph{Verdict:} refuted.
\item[\textbf{Claim 9}] Qwen3.8-27B で SciGym-small 137 件を解かせたとき、反復上限 40 から 80 への平均 STE の改善は 0.05 未満である。~\cite[s1.p3, s1.p4]{duan-2025-measuring}, \cite[s2.p1]{h4duan-2025-scigym}
  \emph{Rationale:} Qwen3.8-27B の small split で、予算を 40 から 80 に倍増しても伸びが判定幅に収まる（飽和している）ことを示し、h2 の当該モデル・split の部分を担う。基準 run は c3 の反復上限 80。
  \emph{Criterion:} \texttt{\detokenize{qwen3.8-27b-small-i40}}.\texttt{\detokenize{trajectory_smape}} $-$ \texttt{\detokenize{qwen3.8-27b-small-i80}}.\texttt{\detokenize{trajectory_smape}} $\leq 0.05$.
  \emph{Prediction:} $[-0.02, 0.04]$ (内部試行では kimi-k3 が 100 反復でも平均 67.8 反復で提出しており、40 を超える予算の使用は部分的。qwen 系は 20 前後で自発提出する。40$\to$80 の差は雑音幅（約 0.03）程度と見る).
  \emph{Observed:} 0.0443245.
  \emph{Verdict:} supported.
\item[\textbf{Claim 10}] Qwen3.8-27B で SciGym-large 213 件を解かせたとき、反復上限 40 から 80 への平均 STE の改善は 0.05 未満である。~\cite[s1.p3, s1.p4]{duan-2025-measuring}, \cite[s2.p1]{h4duan-2025-scigym}
  \emph{Rationale:} Qwen3.8-27B の large split で、予算を 40 から 80 に倍増しても伸びが判定幅に収まる（飽和している）ことを示し、h2 の当該モデル・split の部分を担う。基準 run は c4 の反復上限 80。
  \emph{Criterion:} \texttt{\detokenize{qwen3.8-27b-large-i40}}.\texttt{\detokenize{trajectory_smape}} $-$ \texttt{\detokenize{qwen3.8-27b-large-i80}}.\texttt{\detokenize{trajectory_smape}} $\leq 0.05$.
  \emph{Prediction:} $[-0.02, 0.04]$ (内部試行では kimi-k3 が 100 反復でも平均 67.8 反復で提出しており、40 を超える予算の使用は部分的。qwen 系は 20 前後で自発提出する。40$\to$80 の差は雑音幅（約 0.03）程度と見る).
  \emph{Observed:} 0.07322.
  \emph{Verdict:} refuted.
\item[\textbf{Claim 11}] Nemotron-3-Ultra で SciGym-small 137 件を解かせたとき、反復上限 40 から 80 への平均 STE の改善は 0.05 未満である。~\cite[s1.p3, s1.p4]{duan-2025-measuring}, \cite[s2.p1]{h4duan-2025-scigym}
  \emph{Rationale:} Nemotron-3-Ultra の small split で、予算を 40 から 80 に倍増しても伸びが判定幅に収まる（飽和している）ことを示し、h2 の当該モデル・split の部分を担う。基準 run は c5 の反復上限 80。
  \emph{Criterion:} \texttt{\detokenize{nemotron-3-ultra-small-i40}}.\texttt{\detokenize{trajectory_smape}} $-$ \texttt{\detokenize{nemotron-3-ultra-small-i80}}.\texttt{\detokenize{trajectory_smape}} $\leq 0.05$.
  \emph{Prediction:} $[-0.01, 0.06]$ (内部試行では kimi-k3 が 100 反復でも平均 67.8 反復で提出しており、40 を超える予算の使用は部分的。qwen 系は 20 前後で自発提出する。40$\to$80 の差は雑音幅（約 0.03）程度と見る).
  \emph{Observed:} 0.0264056.
  \emph{Verdict:} supported.
\item[\textbf{Claim 12}] Nemotron-3-Ultra で SciGym-large 213 件を解かせたとき、反復上限 40 から 80 への平均 STE の改善は 0.05 未満である。~\cite[s1.p3, s1.p4]{duan-2025-measuring}, \cite[s2.p1]{h4duan-2025-scigym}
  \emph{Rationale:} Nemotron-3-Ultra の large split で、予算を 40 から 80 に倍増しても伸びが判定幅に収まる（飽和している）ことを示し、h2 の当該モデル・split の部分を担う。基準 run は c6 の反復上限 80。
  \emph{Criterion:} \texttt{\detokenize{nemotron-3-ultra-large-i40}}.\texttt{\detokenize{trajectory_smape}} $-$ \texttt{\detokenize{nemotron-3-ultra-large-i80}}.\texttt{\detokenize{trajectory_smape}} $\leq 0.05$.
  \emph{Prediction:} $[-0.01, 0.06]$ (内部試行では kimi-k3 が 100 反復でも平均 67.8 反復で提出しており、40 を超える予算の使用は部分的。qwen 系は 20 前後で自発提出する。40$\to$80 の差は雑音幅（約 0.03）程度と見る).
  \emph{Observed:} 0.0152157.
  \emph{Verdict:} supported.
\end{enumerate}
\noindent\emph{Assumed, so that the claims together imply Hypothesis 2:}
\begin{itemize}
\item h1 と同じ実行・serving・評価環境の仮定（c7〜c12）
\item 反復上限 40 と 80 の run が同じサーバー設定・同じコミットで実行され、差が予算以外の条件によらないこと（c7〜c12）
\end{itemize}
